\documentclass[10pt,a4paper]{amsart}
\usepackage[margin=19mm]{geometry}
\usepackage{amsmath,amssymb,mathtools}
\usepackage[scr=rsfs,cal=euler]{mathalfa}
\usepackage{xfrac}
\usepackage[unicode,hidelinks,bookmarks=false]{hyperref}
\newcommand{\ie}{i.e.\ }
\newcommand{\cf}{cf.\ }
\newcommand{\resp}{respectively\ }
\newcommand{\Subsection}{\subsection}
\providecommand{\nobreakdash}{\nobreak\hbox{-}}
\setlength{\emergencystretch}{2em}
\multlinegap0pt
\usepackage{bm}
\usepackage{bbm}
\usepackage{dsfont}
\usepackage{xspace}
\usepackage{cancel}
\usepackage{mathtools}
\usepackage{amsthm}
\makeatletter
\@ifundefined{theorem}{\newtheorem{theorem}{Theorem}}{}
\@ifundefined{lemma}{\newtheorem{lemma}[theorem]{Lemma}}{}
\makeatother
\newcommand{\supp}{\operatorname{supp}}
\title[Relative invariant subalgebra rigidity for Thompson's group ]{Relative invariant subalgebra rigidity for Thompson's group $F$}
\author{Tattwamasi Amrutam, Artem Dudko}
\date{Source: arXiv:2606.01268v1 (CC BY 4.0)}
\begin{document}
\maketitle

\noindent\emph{Source and licence.} Relative invariant subalgebra rigidity for Thompson's group $F$. Version arXiv:2606.01268v1 (CC BY 4.0), distributed under the Creative Commons Attribution 4.0 International licence, as recorded in the arXiv licence metadata for this exact version and shown on the abstract page https://arxiv.org/abs/2606.01268v1. Full rights notice: licenses/arxiv-2606.01268-CC-BY-4.0.txt.\par\smallskip

\noindent\textit{Editorial scope.} Counted unit: the Lemma whose source label is \texttt{lem:boundary} in the article (source0.tex lines 261--267, source label \texttt{lem:boundary}), delivered together with the article's own complete original proof of it (source0.tex lines 269--304, one \texttt{proof} environment of 36 lines). No mathematics is written, repaired or re-derived: statement and proof are the article's own text line for line. Required context retained uncounted: none. Every definition, display and label of those lines is retained unchanged. Omitted from this package and not redistributed: 1--260, 268--268, 305--347. Nothing inside a retained excerpt is omitted or altered: every retained body is delivered exactly as the article's lines are written, so no omission receipt of any kind accompanies this package. Citations: the retained excerpts carry no citation command at all, so no bibliography entry of the article is transcribed and no reference is redistributed. Reference adaptation: none was needed and none was made; every reference inside the delivered unit resolves inside it, so no unresolved reference or citation is produced. Printed slips kept literal: no printed word, symbol, hypothesis or number of the counted statement or of its proof was altered. Layout and class: the article's own class is \texttt{amsart} with packages including graphicx, amsmath, amssymb, amsthm, hyperref, xcolor; the wrapper is the lane's pinned \texttt{amsart} editorial wrapper at 10pt on A4, which restates the theorem-like environments the retained text needs and adds the article's own macro definitions that the retained text needs (\texttt{\textbackslash supp}), with extra wrapper packages (bm, bbm, dsfont, xspace, cancel, mathtools). The delivered numbering is the unit's own. Assets: no retained span contains a figure environment, a graphics-inclusion command, a PDF-page inclusion, a table or a TikZ picture; no image file is copied into this package, the package declares no asset dependency and no image file is redistributed with it. Licence: version arXiv:2606.01268v1 of this article is distributed under the Creative Commons Attribution 4.0 International licence, as recorded in the arXiv licence metadata for this exact version and shown on the abstract page https://arxiv.org/abs/2606.01268v1. A full rights notice is delivered at licenses/arxiv-2606.01268-CC-BY-4.0.txt. Characters: every retained excerpt is pure ASCII, so the delivered engine, XeLaTeX with its default Latin Modern text font, reports no missing character.
\begin{lemma}\label{lem:boundary}
Let $x \in L(F)$, and let $\{g_k\} \subset F$ satisfy $\supp(g_k) \to 1$ in the sense that $\supp(g_k) \neq \emptyset$ and $\supp(g_k) \subset (1-\varepsilon_k, 1)$ for some $\varepsilon_k \to 0$. Then $\lambda(g_k)\, x\, \lambda(g_k)^*$ converges in the weak operator topology to an element $\tilde{x} \in L(F)$ whose Fourier coefficients are
\[
\tilde{x}_h = \begin{cases} x_h & \text{if } h'(1) = 1, \\ 0 & \text{if } h'(1) \neq 1, \end{cases} \qquad h \in F.
\]
The analogous statement holds for $\{f_k\} \subset F$ with $\supp(f_k) \to 0$, with $h'(0)$ in place of $h'(1)$.
\end{lemma}

\begin{proof}
Set $x_k := \lambda(g_k)\, x\, \lambda(g_k)^*$. Each $x_k$ has operator norm $\|x\|$, so the sequence lies in the WOT-compact ball $\overline{B}(0, \|x\|) \subset L(F)$, where WOT stands for the weak operator topology.  Let $\{x_{k_i}\}$ be any WOT-convergent subsequence with limit $\tilde{x} \in L(F)$ (the algebra is WOT-closed). For each $h \in F$, the functional $y \mapsto \tau(y\,\lambda(h)^*) = \langle y\delta_e, \delta_h\rangle$ is WOT-continuous, and the trace property gives
\begin{align*}
\tau\!\left(\lambda(g)\, x\, \lambda(g)^* \lambda(h)^*\right) &= \tau\!\left(x\, \lambda(g)^* \lambda(h)^* \lambda(g)\right) = \tau\!\left(x\, \lambda(g^{-1} h g)^*\right) =x_{g^{-1} h g}.
\end{align*}
Hence
\begin{equation}\label{eq:fourier-limit}
\tilde{x}_h \;=\; \lim_{i \to \infty} x_{\, g_{k_i}^{-1} h\, g_{k_i}}.
\end{equation}
\medskip
\noindent\textbf{Case 1: $h'(1) = 1$.} Every element of $F$ is piecewise affine with finitely many dyadic breakpoints and slopes that are integer powers of $2$. On the rightmost piece, the slope equals $h'(1) = 1$; combined with $h(1) = 1$, this forces $h(x) = x$ on some interval $[1 - \delta, 1]$. Once $\varepsilon_{k_i} < \delta$, the support of $g_{k_i}$ lies in a region where $h$ is the identity, so $h$ and $g_{k_i}$ commute, $g_{k_i}^{-1} h g_{k_i} = h$, and equation~\eqref{eq:fourier-limit} gives $\tilde{x}_h = x_h$.

\medskip
\noindent\textbf{Case 2: $h'(1) \neq 1$.} Write $h'(1) = 2^m$ with $m \in \mathbb{Z} \setminus \{0\}$. Choose $\delta > 0$ small enough so that $h$ is affine on $[1 - \delta, 1]$:
\[
h(x) \;=\; 1 - 2^m(1 - x), \qquad x \in [1 - \delta, 1].
\]
On $[1 - \delta, 1)$ the equation $h(x) = x$ reduces to $(1 - x)(1 - 2^m) = 0$, which forces $x = 1$. Thus $h$ has no fixed points in $[1 - \delta, 1)$.
We now claim that if $f \in F$ commutes with $h$ and $\supp(f) \subset (1 - \delta, 1)$, then $f = e$.
Indeed, $fh = hf$ implies $h(\supp(f)) = \supp(f)$. If $f \neq e$, let $a := \inf \supp(f) \in [1 - \delta, 1)$. Since $h$ is an orientation-preserving homeomorphism of $[0,1]$ and preserves $\supp(f)$ as a set, $h(a) = a$, contradicting the absence of fixed points of $h$ in $[1 - \delta, 1)$.

Further, since $\supp(g_{k_i})\to 1$, passing to a subsequence if necessary, we may assume that $g_{k_i}$ are pairwise distinct. Let $s_i := g_{k_i}^{-1}\, h\, g_{k_i}$. For any $i, j$ with $\varepsilon_{k_i}, \varepsilon_{k_j} < \delta$, suppose $s_i = s_j$. Then $f := g_{k_j}\, g_{k_i}^{-1}$ commutes with $h$ and $\supp(f) \subset (1 - \delta, 1)$. Since $fh = hf$, $h$ preserves $\supp(f)$ as a set; if $f \neq e$, then $\inf \supp(f) \in [1 - \delta, 1)$ is a fixed point of $h$, contradicting the choice of $\delta$. Hence $f = e$, which contradicts $g_{k_i}$ being pairwise distinct. Therefore, $s_i$ are pairwise distinct.
% Now suppose $g_{k_i} = g_*$ for $i$ in an infinite set $S \subseteq \mathbb{N}$. For every $i \in S$,
% \[
% \supp(g_*) \;=\; \supp(g_{k_i}) \;\subset\; (1 - \varepsilon_{k_i}, 1).
% \]
% Since $\varepsilon_{k_i} \to 0$ along $S$, for any $\eta > 0$ there exists $i \in S$ with $\varepsilon_{k_i} < \eta$, so $\supp(g_*) \subset (1 - \eta, 1)$. Therefore
% \[
% \supp(g_*) \;\subset\; \bigcap_{\eta > 0} (1 - \eta, 1) \;=\; \emptyset,
% \]
% so $g_* = e$. But then $\supp(g_{k_i}) = \emptyset$ for $i \in S$, contradicting $\supp(g_k) \neq \emptyset$.
% Thus, the $g_{k_i}$ are pairwise distinct for all sufficiently large $i$, and so are the $s_i$ (by the previous paragraph).
Since $x \in L(F) \hookrightarrow \ell^2(F)$, the family $(x_g)_{g \in F}$ is square-summable, hence $x_{s_i} \to 0$. By equation~\eqref{eq:fourier-limit}, $\tilde{x}_h = 0$.

The case $\supp(f_k) \to 0$ is identical with $0$ in place of $1$: choose $\delta$ such that $h$ is affine on $[0, \delta]$, the equation $h(x) = x$ has no solutions in $(0, \delta]$, and the centralizer fact runs verbatim.
\end{proof}

\end{document}
